\documentclass[11pt, letterpaper]{article}

% --- Idioma y codificación ---
\usepackage[spanish]{babel}
\usepackage[utf8]{inputenc}
\usepackage[T1]{fontenc}
\usepackage{lmodern}

% --- Matemáticas ---
\usepackage{amsmath, amssymb, amsthm}

% --- Formato ---
\usepackage[margin=2.5cm]{geometry}
\usepackage{parskip}
\usepackage{booktabs}
\usepackage{hyperref}
\hypersetup{colorlinks=true, linkcolor=blue!70!black, urlcolor=blue!70!black}

% --- Entornos de teoremas ---
\theoremstyle{definition}
\newtheorem{definition}{Definición}[section]
\theoremstyle{plain}
\newtheorem{theorem}{Teorema}[section]
\newtheorem{lemma}[theorem]{Lema}
\newtheorem{corollary}[theorem]{Corolario}
\theoremstyle{remark}
\newtheorem{remark}{Observación}[section]

\title{%
  T03 --- Demostración formal:\\[4pt]
  Inserción amortizada en vector dinámico
}
\author{Curso Linux--C--Rust --- B15 Estructuras de datos en C}
\date{}

\begin{document}
\maketitle
\tableofcontents

% ===================================================================
\section{Preliminar: método del potencial}
% ===================================================================

\begin{definition}[Función de potencial]
Sea $D_i$ el estado de una estructura de datos después de la $i$-ésima
operación. Una \emph{función de potencial} es una función
$\Phi : \text{estados} \to \mathbb{R}$ que satisface:
\begin{enumerate}
  \item $\Phi(D_0) = 0$ \quad (potencial inicial nulo).
  \item $\Phi(D_i) \geq 0$ para todo $i \geq 0$ \quad (potencial no negativo).
\end{enumerate}
\end{definition}

\begin{definition}[Costo amortizado]
El \emph{costo amortizado} de la $i$-ésima operación es:
\begin{equation}\label{eq:amortized}
  \hat{c}_i = c_i + \Phi(D_i) - \Phi(D_{i-1})
\end{equation}
donde $c_i$ es el costo real de la operación.
\end{definition}

La propiedad fundamental es que la suma de costos amortizados acota
superiormente la suma de costos reales:
\begin{equation}\label{eq:fundamental}
  \sum_{i=1}^{n} c_i
  = \sum_{i=1}^{n} \hat{c}_i - \Phi(D_n) + \Phi(D_0)
  \leq \sum_{i=1}^{n} \hat{c}_i
\end{equation}
ya que $\Phi(D_n) \geq 0 = \Phi(D_0)$.

% ===================================================================
\section{Modelo formal}
% ===================================================================

Consideramos un vector dinámico con las siguientes convenciones:

\begin{itemize}
  \item $s$ = \textbf{size}: número de elementos almacenados.
  \item $C$ = \textbf{capacity}: número de posiciones alocadas.
  \item \textbf{Invariante}: $s \leq C$ siempre.
  \item \textbf{Estrategia de crecimiento}: cuando $s = C$ y se ejecuta
    un push, se aloca un nuevo arreglo de tamaño $2C$, se copian los
    $C$ elementos y se inserta el nuevo. Inicialmente $s = 0$, $C = 0$;
    el primer push aloca con capacidad~1.
  \item \textbf{Modelo de costo}: cada escritura (inserción o copia)
    cuesta~1. Un push sin realloc cuesta~1; un push con realloc que
    copia $k$ elementos cuesta $k + 1$.
\end{itemize}

% ===================================================================
\section{Demostración 1: \texorpdfstring{$O(1)$}{O(1)} amortizado
  por método del potencial}
% ===================================================================

\begin{theorem}\label{thm:potential}
La operación \texttt{push} en un vector dinámico con estrategia de
duplicación tiene costo amortizado $O(1)$.
\end{theorem}

Definimos la función de potencial:
\begin{equation}\label{eq:phi}
  \Phi = 2s - C
\end{equation}

\begin{lemma}[$\Phi \geq 0$ siempre]\label{lem:phi-nonneg}
Para todo $i \geq 0$, $\Phi(D_i) \geq 0$.
\end{lemma}

\begin{proof}
Por inducción sobre el número de operaciones.

\textbf{Base.} Estado inicial: $s = 0$, $C = 0$.
$\Phi = 2(0) - 0 = 0 \geq 0$. \checkmark

Tras el primer push (con realloc, $C: 0 \to 1$): $s = 1$, $C = 1$.
$\Phi = 2(1) - 1 = 1 > 0$. \checkmark

\textbf{Paso inductivo.} Supongamos $\Phi \geq 0$ antes de la
$i$-ésima operación. Hay dos casos:

\emph{Caso 1} (sin realloc, $s < C$): Después del push, $s' = s + 1$,
$C' = C$.
\[
  \Phi' = 2(s+1) - C = (2s - C) + 2 = \Phi + 2 \geq 0 + 2 > 0. \quad \checkmark
\]

\emph{Caso 2} (con realloc, $s = C$): Después del push, $s' = s + 1$,
$C' = 2C = 2s$.
\[
  \Phi' = 2(s+1) - 2s = 2 > 0. \quad \checkmark
\]

En ambos casos $\Phi' > 0$, completando la inducción.
\end{proof}

\begin{proof}[Prueba del Teorema~\ref{thm:potential}]
Calculamos el costo amortizado $\hat{c} = c + \Delta\Phi$ en cada caso.

\textbf{Caso 1: push sin realloc} ($s < C$).
\begin{align*}
  c        &= 1 \\
  \Delta\Phi &= \Phi' - \Phi = [2(s+1) - C] - [2s - C] = 2 \\
  \hat{c}  &= 1 + 2 = 3
\end{align*}

\textbf{Caso 2: push con realloc} ($s = C = k$).
Se copian $k$ elementos y se inserta~1:
\begin{align*}
  c          &= k + 1 \\
  \Phi       &= 2k - k = k \\
  \Phi'      &= 2(k+1) - 2k = 2 \\
  \Delta\Phi &= 2 - k \\
  \hat{c}    &= (k+1) + (2-k) = 3
\end{align*}

En ambos casos, $\hat{c} = 3$. Por la propiedad fundamental
(\ref{eq:fundamental}):
\[
  \sum_{i=1}^{n} c_i \leq \sum_{i=1}^{n} \hat{c}_i = 3n.
\]

Por lo tanto, push es $O(1)$ amortizado con constante~3.
\end{proof}

% ===================================================================
\section{Demostración 2: \texorpdfstring{$O(1)$}{O(1)} amortizado
  por método agregado}
% ===================================================================

\begin{theorem}\label{thm:aggregate}
El costo total de $n$ operaciones push en un vector dinámico con
duplicación es menor que $3n$.
\end{theorem}

\begin{proof}
Separamos el costo total:
\[
  T(n) = \underbrace{n}_{\text{inserciones}}
       + \underbrace{\sum_{\text{reallocs}} \text{copias}}_{\text{copias por realloc}}
\]

Partiendo de capacidad~1, los reallocs ocurren en los pushes
$1, 2, 4, 8, \ldots, 2^{\lfloor \log_2 n \rfloor}$. En el realloc del
push $2^j$ se copian $2^j$ elementos. Las copias totales son:
\[
  \sum_{j=0}^{\lfloor \log_2 n \rfloor} 2^j
  = 2^{\lfloor \log_2 n \rfloor + 1} - 1
\]

Acotamos usando $2^{\lfloor \log_2 n \rfloor} \leq n$:
\[
  2^{\lfloor \log_2 n \rfloor + 1} - 1
  = 2 \cdot 2^{\lfloor \log_2 n \rfloor} - 1
  \leq 2n - 1
  < 2n
\]

Por lo tanto:
\[
  T(n) < n + 2n = 3n
\]

El costo amortizado por operación es $T(n)/n < 3 = O(1)$.
\end{proof}

% ===================================================================
\section{Demostración 3: generalización para factor
  \texorpdfstring{$\alpha > 1$}{α > 1}}
% ===================================================================

\begin{theorem}\label{thm:alpha}
Sea $\alpha > 1$ el factor de crecimiento (la capacidad se multiplica
por $\alpha$ en cada realloc). Entonces push es $O(1)$ amortizado, con
cota:
\[
  \hat{c} < \frac{2\alpha - 1}{\alpha - 1}
\]
\end{theorem}

\begin{proof}
Usamos el método agregado. Con capacidad inicial~1, los reallocs ocurren
cuando el size alcanza $\alpha^0, \alpha^1, \alpha^2, \ldots$
(redondeando a enteros). Sea $m = \lfloor \log_\alpha n \rfloor$ el
número de reallocs en $n$ pushes. Las copias totales son:
\[
  \sum_{j=0}^{m} \alpha^j
  = \frac{\alpha^{m+1} - 1}{\alpha - 1}
\]

Como $\alpha^m \leq n$:
\[
  \frac{\alpha^{m+1} - 1}{\alpha - 1}
  \leq \frac{\alpha \cdot n - 1}{\alpha - 1}
  < \frac{\alpha}{\alpha - 1} \cdot n
\]

El costo total es:
\[
  T(n) < n + \frac{\alpha}{\alpha - 1} \cdot n
       = \frac{2\alpha - 1}{\alpha - 1} \cdot n
\]

El costo amortizado por operación:
\[
  \frac{T(n)}{n} < \frac{2\alpha - 1}{\alpha - 1} = O(1)
\]
\end{proof}

\begin{table}[h]
\centering
\begin{tabular}{cc}
\toprule
Factor $\alpha$ & Cota amortizada $\dfrac{2\alpha - 1}{\alpha - 1}$ \\
\midrule
1.5 & 4.0   \\
2   & 3.0   \\
3   & 2.5   \\
4   & 2.33\ldots \\
\bottomrule
\end{tabular}
\caption{La cota mejora conforme $\alpha$ crece: menos reallocs implica
menos copias. Sin embargo, el desperdicio de memoria en el peor caso es
$(\alpha - 1)/\alpha$.}
\end{table}

\begin{remark}
Para $\alpha = 2$ recuperamos la cota $\hat{c} < 3$ del
Teorema~\ref{thm:aggregate}. Para $\alpha = 1.5$ (MSVC
\texttt{std::vector}) la cota es~4. Factores mayores reducen el
amortizado pero desperdician más memoria: $\alpha = 2$ desperdicia
hasta 50\%; $\alpha = 3$ hasta 67\%.
\end{remark}

% ===================================================================
\section{Demostración 4: crecimiento aditivo da
  \texorpdfstring{$\Theta(n)$}{Θ(n)} amortizado}
% ===================================================================

\begin{theorem}\label{thm:additive}
Si la capacidad crece sumando una constante $d > 0$ (en lugar de
multiplicar), el costo total de $n$ pushes es $\Theta(n^2)$ y el costo
amortizado es $\Theta(n)$.
\end{theorem}

\begin{proof}
Con capacidad inicial $d$, los reallocs ocurren en los pushes
$d{+}1,\; 2d{+}1,\; 3d{+}1,\; \ldots$ En los primeros $n$ pushes hay
$m = \bigl\lfloor (n-1)/d \bigr\rfloor$ reallocs. El $j$-ésimo realloc
copia $jd$ elementos.

\textbf{Cota superior.} Las copias totales:
\[
  \sum_{j=1}^{m} jd = d \cdot \frac{m(m+1)}{2}
\]

Como $m \leq n/d$:
\[
  d \cdot \frac{m(m+1)}{2}
  \leq d \cdot \frac{(n/d)(n/d + 1)}{2}
  = \frac{n^2}{2d} + \frac{n}{2}
  = O(n^2)
\]

\textbf{Cota inferior.} Para $n \geq 2d$, tenemos $m \geq n/(2d)$:
\[
  \sum_{j=1}^{m} jd
  \geq d \cdot \frac{(n/(2d))(n/(2d)+1)}{2}
  \geq d \cdot \frac{n^2/(4d^2)}{2}
  = \frac{n^2}{8d}
  = \Omega(n^2)
\]

Combinando: $T(n) = n + \Theta(n^2) = \Theta(n^2)$.

El costo amortizado:
\[
  \frac{T(n)}{n} = \Theta(n)
\]
\end{proof}

\begin{corollary}
La multiplicación por un factor $\alpha > 1$ es la condición necesaria
y suficiente para obtener $O(1)$ amortizado. El crecimiento aditivo,
sin importar cuán grande sea $d$, siempre resulta en costo amortizado
lineal.
\end{corollary}

% ===================================================================
\section{Resumen de resultados}
% ===================================================================

\begin{table}[h]
\centering
\begin{tabular}{llc}
\toprule
\textbf{Resultado} & \textbf{Método} & \textbf{Cota} \\
\midrule
Push con duplicación: $O(1)$ amort.
  & Potencial ($\Phi = 2s - C$)
  & $\hat{c} = 3$ \\
Costo total de $n$ pushes $< 3n$
  & Agregado (serie geométrica)
  & $T(n) < 3n$ \\
Push con factor $\alpha > 1$: $O(1)$ amort.
  & Agregado
  & $\hat{c} < \frac{2\alpha{-}1}{\alpha{-}1}$ \\
Push con crecimiento $+d$: $\Theta(n)$ amort.
  & Agregado (serie aritmética)
  & $T(n) = \Theta(n^2)$ \\
\bottomrule
\end{tabular}
\end{table}

\end{document}
